\documentclass[11pt]{article}
\usepackage{booktabs}
\usepackage[margin=0.8in]{geometry}
\usepackage{graphicx}
\usepackage{url}

\begin{document}
\sloppy

\section*{Problem 05: Concentric-Circle Solver Performance}

This report summarizes the four-by-four-quadrature problem-05 response matrix run
on \texttt{rzwhippet45} on 2026-08-24.  The paper's
\texttt{05/circ\_in\_circ} problem is implemented as challenge-suite case
\texttt{08/circ\_in\_circ}.  The authoritative full-output matrix is
\path{build-rzwhippet-toss_4_x86_64_ib-llvm@19.1.3-release/paper_circ_in_circ_q4x4_response_20260824_152437}.
Commands, logs, response histories, representative renders, status, and validation
tables are archived under
\path{/usr/workspace/tupek2/dev/buckling_solver_paper/results/circ_in_circ/solver_response_20260824_152437}.

\section*{Configuration}

The two-dimensional mesh contains 3047 first-order quadrilateral elements in five
material regions.  Three stiff circular or annular regions use Neo-Hookean density
1000, bulk modulus 1.133333333, and shear modulus 0.2.  Two compliant regions use
density $2.0\times10^{-4}$, bulk modulus $2.266666667\times10^{-7}$, and shear
modulus $3.173333333\times10^{-7}$.  All material integrals use four quadrature
points per coordinate direction, or 16 tensor-product points per quadrilateral.
The bottom boundary is fixed, and a downward top traction is increased in 50 equal
increments to nominal magnitude 0.024.

The matrix uses eight MPI ranks, cold starts, visualization output for every
accepted state, nonlinear tolerance $10^{-11}$, and linear tolerance $10^{-14}$.
Both preconditioners use CG.  Hypre AMG has a 600-iteration cap, and Hypre Jacobi
has a 60000-iteration cap.  The four nonlinear configurations are Newton line
search, PETSc Newton critical-point line search, baseline trust region, and trust
region with safeguarded subspace globalization.  Each run has a 1200-second
allocation.  Both BSR paths are disabled.

This run required a temporary global source change: line 285 of
\path{src/smith/numerics/functional/functional.hpp} sets
\texttt{Functional::Q = 4}.  The build therefore uses four points in each
coordinate direction for every functional, not a case-local override.  Starting
from commit \texttt{9d4024ad9a68b0e1276bfd0b8a4e07e608329f22}, apply the archived
\path{source.patch}, change to
\path{build-rzwhippet-toss_4_x86_64_ib-llvm@19.1.3-release}, run
\texttt{make -j40}, and then run
\path{scripts/run_circ_in_circ_response_matrix.sh}.  The archived
executable has SHA-256
\texttt{4fc62781951e8441e35ae1ca4eccdc2e2e29d1165eeecbe0e401f555d88cef2e}.
The archive retains the source state, patch, executable identity, exact commands,
and postprocessing instructions.

\section*{Completion and Timing}

Three of eight configurations complete all 50 increments, for an overall success
rate of 37.5\%.  Hypre AMG completes two of four configurations, while Hypre Jacobi
completes one of four.  Across the four trust-region configurations the success
rate is 75\%; neither Newton configuration completes with either preconditioner.
Every run has one response state and one visualization cycle for the unloaded state
and each accepted increment.

\begin{table}[h]
  \centering
  \small
  \caption{Problem-05 completion and reported solver time.  Timeout entries use launcher elapsed time.}
  \begin{tabular}{@{}lll@{}}
    \toprule
    Solver & Hypre AMG & Hypre Jacobi \\
    \midrule
    Newton-LS & $23/50$ steps (timeout, $1201\,\mathrm{s}$) & $23/50$ steps (timeout, $1202\,\mathrm{s}$) \\
    Newton-CP & $3/50$ steps (failure, $0.671\,\mathrm{s}$) & $3/50$ steps (failure, $1.158\,\mathrm{s}$) \\
    Trust-region & $50/50$ steps, $18.005\,\mathrm{s}$ & $34/50$ steps (failure, $76.347\,\mathrm{s}$) \\
    Trust-region + subspace & $50/50$ steps, $20.780\,\mathrm{s}$ & $50/50$ steps, $40.199\,\mathrm{s}$ \\
    \bottomrule
  \end{tabular}
\end{table}

The baseline trust-region/AMG configuration is the fastest completed run.  The
subspace safeguard adds $2.775\,\mathrm{s}$ with AMG.  With Jacobi, however,
the safeguard changes the outcome from failure at increment 35 to full completion,
at about twice the wall time of the completed AMG configurations.

\begin{figure}[h]
  \centering
  \includegraphics[width=0.86\textwidth]{/usr/workspace/tupek2/dev/buckling_solver_paper/figures/circ_in_circ_solver_timing.png}
  \caption{Log-scale reported solver time and completion status for the eight problem-05 configurations; timeout entries use launcher elapsed time.}
\end{figure}

\section*{Response Branches}

All partial histories coincide with the completed reference branch at the precision
written to CSV.  Newton-CP accepts three increments and stops at applied traction
0.00144 and average top displacement $-0.0159179$.  Newton-LS accepts 23 increments
and reaches traction 0.01104 and displacement $-0.171289$ before spending the
remainder of its allocation on increment 24.  Baseline trust region with Jacobi
accepts 34 increments and reaches traction 0.01632 and displacement $-0.470531$
before failing on increment 35.

The completed solutions traverse the sharp buckling transition between increments
34 and 35: average top displacement changes from $-0.470531$ at traction 0.01632 to
$-0.595089$ at traction 0.0168.  All three completed histories reach traction 0.024
and final displacement between $-0.615707$ and $-0.615708$.  Their largest
pointwise displacement spread is $8.0\times10^{-6}$, so they are indistinguishable
at plotting resolution.

\begin{figure}[h]
  \centering
  \begin{minipage}{0.49\textwidth}
    \centering
    \includegraphics[width=\linewidth]{/usr/workspace/tupek2/dev/buckling_solver_paper/figures/circ_in_circ_solver_response_HypreAMG.png}
  \end{minipage}
  \begin{minipage}{0.49\textwidth}
    \centering
    \includegraphics[width=\linewidth]{/usr/workspace/tupek2/dev/buckling_solver_paper/figures/circ_in_circ_solver_response_HypreJacobi.png}
  \end{minipage}
  \caption{Applied-traction versus average top-displacement histories.  Coincident curves overlay one another.}
\end{figure}

\section*{Solver Diagnostics}

Both Newton-CP runs report PETSc convergence reason $-4$ on increment 4.  Baseline
trust region with Jacobi reaches its 10000-iteration nonlinear limit on increment
35.  This is a nonlinear stagnation rather than an inner-CG failure: the run
records neither an indefinite-CG diagnostic nor a hit of the 60000-iteration
linear cap.  Jacobi-preconditioned nonlinear work rises sharply approaching the
branch transition, requiring 1638, 2630, and 5562 iterations on increments 32,
33, and 34 before the failed increment reaches the cap with residual
$6.05074\times10^{-5}$.  The safeguarded subspace configuration crosses increment
35 in 4698 iterations by augmenting the baseline Steihaug--Toint/dogleg step with
the Cauchy direction, recent accepted nonlinear steps, and recycled leftmost-
curvature information.  Across accepted increments, baseline trust region/AMG uses 929 nonlinear
iterations, safeguarded trust region/AMG uses 1059, and safeguarded trust
region/Jacobi uses 8133.  The two Newton-LS logs contain 7911 and 82875
indefinite-CG diagnostics for AMG and Jacobi, respectively, before timing out.
In contrast, the safeguarded trust-region method completes the prescribed loading
with both preconditioners.  This is the principal solver-robustness separation in
the matrix.

\section*{Deformation Comparison}

Each row below samples up to six states evenly over that run's accepted history;
columns therefore represent progress within a row rather than identical load
levels across rows.  Newton-CP has only four available states, so intermediate
panels are repeated to retain the common layout.  The AMG and Jacobi comparisons
show the same accepted deformation path.  Only the completed trust-region rows
reach the strongly collapsed final configuration.  Colors identify the five
material regions, mesh edges remain visible, and the deformed geometry uses unit
displacement warp.

\begin{figure}[h]
  \centering
  \includegraphics[width=\textwidth]{/usr/workspace/tupek2/dev/buckling_solver_paper/figures/challenge_snapshots/challenge_circ_in_circ_solver_comparison_HypreAMG.png}
  \caption{Accepted-state deformation comparison with Hypre AMG.}
\end{figure}

\begin{figure}[h]
  \centering
  \includegraphics[width=\textwidth]{/usr/workspace/tupek2/dev/buckling_solver_paper/figures/challenge_snapshots/challenge_circ_in_circ_solver_comparison_HypreJacobi.png}
  \caption{Accepted-state deformation comparison with Hypre Jacobi.}
\end{figure}

\section*{Third-Medium Deformation Localization}

A 17-by-17 reference-coordinate sampling of each quadrilateral finds the first
negative deformation Jacobian on cycle 14 for all three completed branches.  The
material-attribute-resolved check shows that every such element belongs to one of
the two highly compliant third-medium regions.  Attributes 1, 3, and 5 are the
three stiff regions; attributes 2 and 4 are the compliant regions.

\begin{table}[h]
  \centering
  \small
  \caption{Material localization of nonpositive sampled deformation Jacobians.}
  \begin{tabular}{@{}lrrrr@{}}
    \toprule
    Completed branch & First cycle & Cells at first cycle & Minimum attribute & Final counts by attribute \\
    \midrule
    Trust-region/AMG & 14 & 8 & 2 & $2{:}85,\ 4{:}3$ \\
    Subspace trust-region/AMG & 14 & 8 & 2 & $2{:}85,\ 4{:}3$ \\
    Subspace trust-region/Jacobi & 14 & 8 & 2 & $2{:}85,\ 4{:}3$ \\
    \bottomrule
  \end{tabular}
\end{table}

For baseline trust region/AMG, the sampled minimum changes from
$9.4414\times10^{-4}$ on cycle 13 to $-3.6749\times10^{-4}$ on cycle 14.  Its
minimum over the run is $-1.1641912333$ on cycle 45 in attribute 2.  At cycle 50,
85 inverted quadrilaterals are in attribute 2 and three are in attribute 4; none
are in a stiff region.  The two safeguarded branches give the same localization
and counts, with differences only at roundoff level.

The third medium is intentionally sacrificial and six orders of magnitude softer
than the stiff bodies.  The localized inversion is not visible in the unit-warp
geometry and does not alter the coincident response histories of the three
completed branches.  For the purpose of this paper, the calculation therefore
supports the intended contact-like load-transfer and nonlinear-solver robustness
comparison.  It is not used to claim an accurate stress or deformation field
inside the third medium, formal contact convergence, or physical admissibility of
that fictitious material.  The attribute-resolved tables are retained as
\path{deformation_validity_by_attribute.tsv} with each completed branch.

\end{document}
